%------------------------------------------------
%--Nombre de archivo código, Semestre, año--------- 
%---------------"F702-P2015"------Primer Semestre---
%---------------"F702-S2015"------Segundo Semestre---
%--------------------------------------------------
\documentclass{article}

\usepackage{amsmath}
\usepackage{amsfonts}
\usepackage{amssymb}
\usepackage{graphicx}
\usepackage{fancyhdr}
\usepackage[utf8]{inputenc} 
%\pagestyle{fancyplain}
\usepackage{hyperref}
\usepackage{enumerate}

\setcounter{section}{0}

\oddsidemargin=0.5cm
\setlength{\textwidth}{6.5in}
\textheight 9.0in 
\addtolength{\voffset}{-0.8in}
\lhead{{\bf ECFM USAC} }

% Macros para automatizar contenidos de las unidades
\newcommand{\unidades}[4]{%
\begin{description}
	\item[Descripción:] #1
	\item[Duración:] #2
	\item[Metodología:] #3	
	\item[Evaluación:] #4
\end{description}	
}

\begin{document}
%---------------------------------------------------------------------
\noindent{\large Universidad de San Carlos de Guatemala\\
Escuela de Ciencias Físicas y Matemáticas\\
Programas de Licenciatura}
\bigskip

\bigskip
\centerline{\Large Programa Laboratorio de Electrónica Digital}
%---------------------------------------------------------------------
\section{Descripción del Curso}
\begin{tabular}{ll}
{\bf Nombre:} Laboratorio de electrónica digital & {\bf Código:} F604\\
{\bf Prerrequisitos:} F503  & {\bf Créditos:} 5\\
{\bf Profesor:} Erick Díaz & {\bf Semestre:} Segundo, 2026\\
\end{tabular}
\bigskip

Es un curso introductorio que se ocupa del diseño de circuitos electrónicos digitales. Los circuitos digitales
se emplean en el diseño y construcción de sistemas como computadoras digitales, comunicación de datos, grabación digital, instrumentación
física y muchas otras aplicaciones que requieren hardware digital. Se finaliza el curso con una introducción a programación en lenguaje C.
%-------------------------------------------------------------------
\section{Competencias}
\subsection{Competencias generales} % AQUI SE UTILIZA ENUMERATE
\begin{enumerate}[2.{1}.1]
\item Utilizar o elaborar programas o sistemas de computación para el procesamiento de información, cálculo numérico, simulación de procesos físicos o control de experimentos.

\item Construir modelos simplificados que describan una situación compleja, identificando sus elementos esenciales y efectuando las aproximaciones necesarias.

\item Verificar y evaluar el ajuste de modelos a la realidad, identificando su dominio de validez.

\item Aplicar el conocimiento teórico de la física en la realización e interpretación de experimentos.

\item Desarrollar argumentaciones válidas en el ámbito de la física, identificando hipótesis y conclusiones.

\item Estimar el orden de magnitud de cantidades mensurables para interpretar fenómenos diversos.


 

\end{enumerate}


\subsection{Competencias específicas}
\begin{enumerate}[\hspace{3ex}a)]
\item Demostrar destrezas experimentales y uso de métodos adecuados de trabajo en el laboratorio.
\item Describir y explicar fenómenos naturales y procesos tecnológicos en términos de conceptos, principios y teorías físicas.
\item Percibir las analogías entre situaciones aparentemente diversas, utilizando soluciones conocidas en la resolución de problemas nuevos.
\item Demostrar hábitos de trabajo necesarios para el desarrollo de la profesión tales como el trabajo en equipo, el rigor científico, el auto-aprendizaje y la persistencia. 

\end{enumerate}
%-----------------------------------------------------------
\section{Unidades}


\subsection{Sistemas binarios}

\unidades{Números binarios y decimales, Conversión decimal-binario, Números hexadecimales, Códigos binarios, Lógica binaria.}{4 periodos de 50 minutos}{Los períodos de clase son magistrales, con la solución de algunos ejercicios guías, para que el estudiante demuestre su aprendizaje con la resolución de los ejercicios propuestos.}{Se evaluará por medio de examen teórico}

\subsection{Álgebra booleana y compuertas lógicas}
\unidades{Definiciones básicas, Propiedades del álgebra booleana, Teoremas del álgebra booleana, Funciones booleana, Formas canónicas y estandar, Operadores NOT, Operación OR, Operación AND, Otras operaciones básicas, Mapas de Karnough}{12 períodos de 50 minutos}{Los períodos de clase son magistrales, con la solución de algunos ejercicios guías, para que el estudiante demuestre su aprendizaje con la resolución de los ejercicios propuestos y practicas de laboratorio planificadas.}{Se evaluará por medio de examen teórico e informes de las prácticas de laboratorio realizadas.}

\subsection{Lógica Combinacional}
\unidades{Procedimiento de Diseño, Sumadores, Conversión de códigos, Multiplexores.}{8 períodos de 50 minutos}{Los períodos de clase son magistrales, con la solución de algunos ejercicios guías, para que el estudiante demuestre su aprendizaje con la resolución de los ejercicios propuestos y desarrollo de mini-proyectos.}{Se evaluará por medio de examen teórico e informes de mini-proyectos.}

\subsection{Lógica Secuencial}
\unidades{Flip-Flops, Procedimiento de Diseño, Sincronización, Contadores, Registros, Secuencias temporizadas.}{8 períodos de 50 minutos}{Los períodos de clase son magistrales, con la solución de algunos ejercicios guías, para que el estudiante demuestre su aprendizaje con la resolución de los ejercicios propuestos y desarrollo de mini-proyectos.}{Se evaluará por medio de examen teórico e informes de mini-proyectos.}


\subsection{Introducción a los microprocesadores}
\unidades{Compomentes de un microprosesador. Tipos de microprocesadores. Lenguaje máquina. Sistemas operativos}
{4 períodos de 50 minutos}{Los períodos de clase son magistrales, con la solución de algunos ejercicios guías, para que el estudiante demuestre su aprendizaje con la resolución de los ejercicios propuestos.}{Se evaluará por medio de una tarea.}

\subsection{Programación en C}
\unidades{Compilador. Tipos de variables. Estructuras condicionales. Estructuras de bucle. Funciones. Manejo de memoria.}
{24 períodos de 50 minutos}{Los períodos de clase son magistrales, con la solución de algunos ejercicios guías, para que el estudiante demuestre su aprendizaje con la resolución de los ejercicios propuestos y practicas de laboratorio planificadas.}{Se evaluará por medio de examen teórico y tareas.}

\subsection{Introducción a FPGA}
\unidades{Conceptos básicos de FPGA. Flujo de diseño en FPGA. Hardware Description Language (HDL). Síntesis y programación. Aplicaciones de FPGA.}{6 períodos de 50 minutos}{Los períodos de clase son magistrales con demostraciones de herramientas CAD, combinados con prácticas de laboratorio guiadas para la implementación de circuitos digitales en FPGA.}{Se evaluará por medio de informes de prácticas de laboratorio.}

%----------------------------------------------------------------

\section{Evaluación del curso}
Los porcentajes asignados a cada uno de los elementos de la evaluación están de acuerdo con el Reglamento General de Evaluación y Promoción del Estudiante de la Universidad de San Carlos de Guatemala 
\bigskip

\begin{tabular}{lr}
	 Examenes parciales & 40 puntos\\
	 Practicas Laboratorio & 20 puntos\\
	 Tareas & 15 puntos\\
	Evaluación final & 25 puntos\\\hline
	Total & 100 puntos\\
\end{tabular}
%--El reglamento establece que la zona debe estar entre el 70% y el 85% 
%--de la nota del curso
%-------------------------------------------------------------------

\section{Bibliografía}
\begin{enumerate}
\item M. Morris Mano, Diseño Digital, Pearson, 2003, Tercera edición, México.
\item P. J. Deitel, H. M. Deitel, C++ Como Programar, Pearson, 2008, Sexta Edición, México.

\end{enumerate}
\bigskip
\flushright{http://ecfm.usac.edu.gt/programas}
\end{document}

---------------------------------------------------------------

COMPETENCIAS 

COMPETENCIAS GENERALES
En el programa usted debe poner todas las competencias que aparecen en el Diseño Curricular que aplican al curso, en orden de importancia en el curso, no es necesario que incluya todas
---------------------------------------------------------------------
COMPETENCIAS PRESENTES EN EL DISEÑO CURRICULAR
Los estudiantes al graduarse de la licenciatura en Física deben poseer las
siguientes competencias:

a) Plantear, analizar y resolver problemas físicos, tanto teóricos como
experimentales, mediante la utilización de métodos analíticos,
experimentales o numéricos.

b) Utilizar o elaborar programas o sistemas de computación para el
procesamiento de información, cálculo numérico, simulación de
procesos físicos o control de experimentos.

c) Construir modelos simplificados que describan una situación
compleja, identificando sus elementos esenciales y efectuando las
aproximaciones necesarias.

d) Verificar y evaluar el ajuste de modelos a la realidad, identificando su
dominio de validez.

e) Aplicar el conocimiento teórico de la física en la realización e
interpretación de experimentos.

f) Demostrar una comprensión profunda de los conceptos y principios
fundamentales, tanto de la física clásica como de la física moderna.

g) Describir y explicar fenómenos naturales y procesos tecnológicos en
términos de conceptos, principios y teorías físicas.

h) Desarrollar argumentaciones válidas en el ámbito de la física,
identificando hipótesis y conclusiones.

i) Sintetizar soluciones particulares, extendiéndolas hacia principios,
leyes o teorías más generales.

j) Percibir las analogías entre situaciones aparentemente diversas,
utilizando soluciones conocidas en la resolución de problemas
nuevos.

k) Estimar el orden de magnitud de cantidades mensurables para
interpretar fenómenos diversos.

l) Demostrar destrezas experimentales y uso de métodos adecuados
de trabajo en el laboratorio.

m) Participar en actividades profesionales relacionadas con tecnologías
de alto nivel, sea en el laboratorio o en la industria.

n) Participar en asesorías y elaboración de propuestas en ciencia y
tecnología en temas con impacto económico y social en el ámbito
nacional.

o) Actuar con responsabilidad y ética profesional, manifestando
conciencia social de solidaridad, justicia, y respeto por el ambiente.

p) Demostrar hábitos de trabajo necesarios para el desarrollo de la
profesión tales como el trabajo en equipo, el rigor científico, el auto-
aprendizaje y la persistencia.

q) Buscar, interpretar y utilizar información científica.

r) Comunicar conceptos y resultados científicos en lenguaje oral y
escrito ante sus pares, y en situaciones de enseñanza y de
divulgación.

s) Participar en la elaboración y desarrollo de proyectos de
investigación en física interdisciplinarios.

t) Demostrar disposición para enfrentar nuevos problemas en otros
campos, utilizando sus habilidades y conocimientos específicos.

u) Conocer y comprender el desarrollo conceptual de la física en
términos históricos y epistemológicos.

v) Conocer los aspectos relevantes del proceso de enseñanza-
aprendizaje de la física, demostrando disposición para colaborar en
la formación de científicos.
-----------------------------------------------------------------
COMPETENCIAS ESPECIFICAS
Usted puede agregar competencias específicas del curso que esten contenidas en las Competencias Generales
------------------------------------------------------------------------




